\documentclass[12pt]{article}
\usepackage[utf8]{inputenc}
\usepackage[T1]{fontenc}
\usepackage{hyperref}
\title{Moiso Language Resources for K-12 Curriculum Development}
\author{Compiled by Claude Sonnet 5.5 (claude-sonnet-5-5) for Andr\'as Kornai}
\date{}
\begin{document}
\maketitle

\section{Introduction}
Moiso (ISO 639-3 wlr; Glottocode wail1242) is an Austronesian (Oceanic) language of Espiritu Santo, Vanuatu. Sources disagree on its name and size. The ISO code wlr is also listed under the names Wailapa and Ale. Wikipedia places Moiso in Moriuli village in central Santo and gives about 100 speakers, citing Tryon (2010). Joshua Project gives a population of roughly 300. Very little documentation turned up in the searches. No school materials were found.

\section{Primary Educational Resources}
No primers, textbooks or school programmes in Moiso were found in the searches. A curriculum would have to be built from scratch with the community.

\section{Linguistic and Reference Materials}
\begin{itemize}
\item Wikipedia, Moiso language: \url{https://en.wikipedia.org/wiki/Moiso_language}
\item Wikipedia, Wailapa language, which covers the same ISO code under the names Wailapa, Moiso and Ale and notes a possible link to a proposed Southwest Santo language: \url{https://en.wikipedia.org/wiki/Wailapa_language}
\end{itemize}
No grammar or dictionary was found in the search results.

\section{Religious and Community Resources}
\begin{itemize}
\item Joshua Project people-group profile (Wailapa in Vanuatu; main language listed as Moiso): \url{https://joshuaproject.net/people_groups/15823/NH}
\end{itemize}

\section{Digital and Audio Resources}
\begin{itemize}
\item Global Recordings Network page for code wlr (listed as Moiso; other names Ale, Biliru, Wailapa). I did not check which recordings it holds: \url{https://globalrecordings.net/en/language/wlr}
\end{itemize}

\section{Curriculum Development Context}
Documentation is thin and the sources conflict on speaker numbers and on the relation between Moiso and Wailapa. Any curriculum would need fieldwork with the Moriuli community. It would probably rely on Bislama and English as the school languages, with Moiso used for oral and cultural content.

\end{document}